% =================================================================
% BÖLÜM 3 — SENSÖR FÜZYONU VE UÇUŞ ALGORİTMASI
% =================================================================
\section{Sensör Füzyonu ve Uçuş Algoritması}

% ---------------------------------------------------------------
\begin{frame}{IMU Veri İşleme Hattı ve Mahony AHRS}{Yönelim kestirimi (Roll / Pitch / Yaw)}
  \centering
  \includegraphics[width=\textwidth,height=0.46\textheight,keepaspectratio]{imu_pipeline.png}
  \par\vspace{2pt}
  \begin{columns}[T]
    \begin{column}{0.36\textwidth}
      \begin{block}{\tiny Neden füzyon?}
        \tiny İvmeölçer titreşimde \textbf{gürültülü}; jiroskop zamanla \textbf{kayar}. İkisinin güçlü yönü birleştirilir.
      \end{block}
    \end{column}
    \begin{column}{0.64\textwidth}
      \begin{exampleblock}{\tiny Mahony AHRS — Nasıl Çalışır?}
        \tiny
        \textbf{AHRS} = Attitude \& Heading Ref. System (Roll/Pitch/Yaw)\\[1pt]
        \textbf{1.} Jiro $\rightarrow$ açısal hız $\rightarrow$ quaternion güncelle\quad
        \textbf{2.} İvmeölçer $\rightarrow$ yerçekimi ref. $\rightarrow$ hata hesapla\quad
        \textbf{3.} PI $\rightarrow$ düzelt\\[1pt]
        \textit{Jiro hız bilgisi verir; ivmeölçer ``kayıyorsun, düzelt'' der.}\\[1pt]
        \textbf{$k_p$} = düzeltme hızı ($k_p{\uparrow}$ titreşime duyarlı, $k_p{\downarrow}$ yavaş)\enspace
        \textbf{$k_i$} = bias düzeltmesi\\
        SIT'te statik açı salınımına bakılarak ayarlandı.\enspace Yakınsama: \textbf{2--3 s}.
      \end{exampleblock}
    \end{column}
  \end{columns}
\end{frame}

% ---------------------------------------------------------------
\begin{frame}{Yükseklik Kalman Filtresi}{Baro + IMU füzyonu ile gürültüden arınmış irtifa}
  \centering
  \includegraphics[width=\textwidth,height=0.44\textheight,keepaspectratio]{kalman.png}
  \par\vspace{2pt}
  \begin{columns}[T]
    \begin{column}{0.40\textwidth}
      \begin{block}{\tiny 3-durumlu kestirim}
        \tiny \textbf{[irtifa,\ hız,\ ivme-bias]}\\
        Baro: yavaş ama mutlak ref.\\
        IMU: hızlı ama kayma var\\
        Kalman $\rightarrow$ güvenilir \textbf{hız} $\rightarrow$ doğru apogee
      \end{block}
      \vspace{1pt}
      \begin{alertblock}{\tiny Q ve R}
        \tiny
        \textbf{R} = ölçüm gürültüsü (baro varyansı ölçüldü)\\
        \textbf{Q} = süreç gürültüsü (model sapması)\\
        Q${\downarrow}$ sensöre güven / Q${\uparrow}$ modele güven
      \end{alertblock}
    \end{column}
    \begin{column}{0.60\textwidth}
      \begin{exampleblock}{\tiny Nasıl Çalışır? — Predict \& Update}
        \tiny
        \textbf{Predict:} fizik modeli ile sonraki durumu tahmin et\\
        \quad irtifa $+{=}$ hız$\times$dt,\enspace hız $+{=}$ ivme$\times$dt\\[2pt]
        \textbf{Update:} baro ölçümü gelince tahmini düzelt\\[2pt]
        \textbf{R} = baroya ne kadar güvenilsin?\\
        \quad R${\uparrow}$ $\rightarrow$ az ağırlık\enspace/\enspace R${\downarrow}$ $\rightarrow$ çok ağırlık\\
        \quad Baro gürültülü $\Rightarrow$ R yüksek tutuldu.
      \end{exampleblock}
    \end{column}
  \end{columns}
\end{frame}

% ---------------------------------------------------------------
\begin{frame}{Uçuş Durum Makinesi — 7 Faz}{If-else spagettisi yerine izlenebilir FSM}
  \begin{columns}[c]
    \begin{column}{0.58\textwidth}
      \centering
      \includegraphics[width=\linewidth,height=0.64\textheight,keepaspectratio]{flight_fsm.png}
    \end{column}
    \begin{column}{0.42\textwidth}
      \scriptsize
      \textbf{IDLE} $\rightarrow$ \textbf{BOOST} $\rightarrow$ \textbf{COAST} $\rightarrow$
      \textbf{APOGEE} $\rightarrow$ \textbf{DROGUE} $\rightarrow$ \textbf{MAIN} $\rightarrow$ \textbf{LANDED}
      \vspace{4pt}
      \begin{alertblock}{\footnotesize Çift güvenlikli apogee}
        \scriptsize
        \begin{itemize}
          \item \textbf{Hız bazlı:} Kalman dikey hızı $<$ 0
          \item \textbf{Tilt bazlı:} Pitch açısı $>$ 70$^\circ$
        \end{itemize}
        Biri tetiklenince Drogue paraşütü açılır.
      \end{alertblock}
      \begin{itemize}
        \item BOOST: ivme $>$ 45 m/s$^2$
        \item MAIN: irtifa $<$ 300 m
        \item LANDED: hız $<$ 2 m/s
      \end{itemize}
    \end{column}
  \end{columns}
\end{frame}

% ---------------------------------------------------------------
\begin{frame}{Güvenilirlik Mekanizmaları}{Kilitlenmeye ve sensör arızasına karşı koruma}
  \begin{columns}[T]
    \begin{column}{0.5\textwidth}
      \begin{block}{Donanımsal / yazılımsal kalkanlar}
        \scriptsize
        \begin{itemize}\setlength{\itemsep}{3pt}
          \item \textbf{IWDG watchdog} ($\sim$2 s)\\
            \tiny Donanım sayacı geri sayar; yazılım 2 sn içinde sıfırlamazsa CPU reset atar.
            Baro görevi her 100 ms'de bir besler — donma olursa otomatik kurtarma.
          \item \textbf{Stack overflow koruması}\\
            \tiny Her FreeRTOS task'ının ayrı stack'i var. \texttt{configCHECK\_FOR\_STACK\_OVERFLOW=2}:
            FreeRTOS her context-switch'te stack dolgusunu kontrol eder; taşma varsa hook çağrılır.
          \item \textbf{Fatal handler} $\rightarrow$ \texttt{NVIC\_SystemReset()}\\
            \tiny Hata durumunda sonsuz döngüde kalmak yerine sistemi yeniden başlatır.
            IWDG + reset = sistem kendi kendine toparlanabilir.
        \end{itemize}
      \end{block}
    \end{column}
    \begin{column}{0.5\textwidth}
      \begin{exampleblock}{Baro-only ``degraded'' mod}
        \footnotesize IMU başlamazsa sistem \textbf{kaybolmaz, uçmaya devam eder}:
        \scriptsize
        \begin{tabular}{@{}lll@{}}
          \toprule
          \textbf{Özellik} & \textbf{IMU var} & \textbf{IMU yok} \\
          \midrule
          BOOST tespiti & ivme & baro hızı \\
          Kalman girdisi & dikey ivme & 0 \\
          Tilt apogee & aktif & pasif \\
          Hız apogee & aktif & \textbf{aktif} \\
          \bottomrule
        \end{tabular}
      \end{exampleblock}
      \centering\scriptsize
      IMU sürücüsü 3 deneme başarısız olursa görev silinir, sistem baro moduyla devam eder.
    \end{column}
  \end{columns}
\end{frame}
